\section{Compressed Krylov and CKCI}
\label{sec:ck-ckci}

We write the many-body wavefunction in a determinant basis as
\begin{equation}
|\Psi\rangle = \sum_I c_I |D_I\rangle,
\end{equation}
where $|D_I\rangle$ denotes a Slater determinant (or, more generally, a
configuration basis state). In the compressed Krylov (CK) algorithm, we build a
sequence of sparse Krylov vectors
\begin{equation}
\mathcal{Q}_m = \mathrm{span}\left\{ |\Psi_0\rangle, |\Psi_1\rangle, \dots, |\Psi_{m-1}\rangle \right\},
\end{equation}
starting from an initial sparse seed $|\Psi_0\rangle$, typically a
Hartree--Fock determinant or another guiding configuration.

At Krylov step $k$, the projected update is generated by Hamiltonian action on
the current sparse vector,
\begin{equation}
|\widetilde{v}_{k+1}\rangle = H |\Psi_k\rangle.
\end{equation}
To remove components already represented in the current Krylov basis, we form
an orthogonalized residual
\begin{equation}
|r_{k+1}\rangle = \left( I - P_k \right)|\widetilde{v}_{k+1}\rangle,
\end{equation}
where $P_k$ denotes projection onto the span of the previously constructed
Krylov vectors. In a nonorthogonal Krylov basis this projector may be written
formally as
\begin{equation}
P_k = \sum_{i,j=0}^{k} |\Psi_i\rangle (S_k^{-1})_{ij} \langle \Psi_j |,
\qquad
(S_k)_{ij} = \langle \Psi_i | \Psi_j \rangle.
\end{equation}
Because the stored Krylov vectors are compressed and generally nonorthogonal,
the corresponding projected eigenvalue problem is written in the generalized
Ritz form
\begin{equation}
H^{(m)} \mathbf{u} = E\, S^{(m)} \mathbf{u},
\end{equation}
with
\begin{equation}
H^{(m)}_{ij} = \langle \Psi_i | H | \Psi_j \rangle,
\qquad
S^{(m)}_{ij} = \langle \Psi_i | \Psi_j \rangle.
\end{equation}
The associated Ritz state has the CI-like expansion
\begin{equation}
|\Phi^{(m)}\rangle = \sum_{k=0}^{m-1} u_k |\Psi_k\rangle.
\end{equation}
The lowest Ritz root gives the current CK energy estimate, and the associated
Ritz vector defines the best approximation within the compressed Krylov space.

Compression enters at two levels. First, only the largest-amplitude components
needed for propagation are retained under a propagation budget
$K_{\mathrm{prop}}$. Second, the next stored Krylov vector is compressed to a
storage budget $K_{\mathrm{store}}$. Formally, if $\mathcal{C}_{K}$ denotes the
operation that retains the $K$ largest coefficients by magnitude, the stored
next vector can be written schematically as
\begin{equation}
|\Psi_{k+1}\rangle =
\frac{\mathcal{C}_{K_{\mathrm{store}}}\!\left(|r_{k+1}\rangle\right)}
{\left\|\mathcal{C}_{K_{\mathrm{store}}}\!\left(|r_{k+1}\rangle\right)\right\|}.
\end{equation}
In the thick-restart form of CK, several low-lying Ritz vectors are retained at
the end of one cycle and used as restart seeds for the next cycle, rather than
restarting from a single vector.

The CKCI construction uses the determinant support discovered by CK as the
selected CI space. If
\begin{equation}
\Omega_{\mathrm{CK}} = \bigcup_{k=0}^{m-1} \mathrm{supp}(\Psi_k)
\end{equation}
is the union of determinant supports appearing in the CK basis, then CKCI-0 is
defined by exact diagonalization of the Hamiltonian projected onto this support,
\begin{equation}
|\Psi_{\mathrm{CKCI\mbox{-}0}}\rangle
= \sum_{I \in \Omega_{\mathrm{CK}}} c_I^{(0)} |D_I\rangle,
\end{equation}
with coefficients obtained from
\begin{equation}
P_{\Omega_{\mathrm{CK}}} H P_{\Omega_{\mathrm{CK}}}
|\Psi_{\mathrm{CKCI\mbox{-}0}}\rangle
= E_{\mathrm{CKCI\mbox{-}0}}
|\Psi_{\mathrm{CKCI\mbox{-}0}}\rangle.
\end{equation}

Higher-order CKCI refinements are obtained by enlarging the determinant support
from the dominant amplitudes of the previous exactized solution. If
$\Omega_n$ denotes the determinant space for CKCI-$n$, and
$\mathcal{P}_{N_{\mathrm{top}}}(\Psi_{\mathrm{CKCI\mbox{-}n}})$ denotes the set
of the $N_{\mathrm{top}}$ determinants with the largest $|c_I^{(n)}|$, then a
single shell-growth refinement may be written schematically as
\begin{equation}
\Omega_{n+1}
=
\Omega_n \cup \mathcal{N}\!\left(
\mathcal{P}_{N_{\mathrm{top}}}
(\Psi_{\mathrm{CKCI\mbox{-}n}})
\right),
\end{equation}
where $\mathcal{N}$ denotes the set of configurations connected by one
application of $H$. CKCI-1 and CKCI-2 are then obtained by exact
diagonalization in $\Omega_1$ and $\Omega_2$, respectively. In this way, CK
acts as a compressed Krylov discovery stage, while CKCI provides an exactized
selected-CI refinement on the determinant support uncovered by the Krylov
dynamics.
